\title{2024 年高考数学 上海卷}
\date{}
\maketitle

\section*{一、填空题}
本大题共有12题，满分54分.其中第1-6题每题4分，第7-12题每题满分5分

\begin{question}[points=4]
设全集$U=\{1,2,3,4,5\}$，集合$A=\{2,4\}$，则$\overline{A}=$.
\fillin[{$\{1,3,5\}$}]
\begin{solution}
\textbf{答案：}$\{1,3,5\}$

\textbf{分析：}根据补集的定义可求$\overline{A}$.

\textbf{详解：}由题设有$\overline{A}=\{1,3,5\}$，故答案为：$\{1,3,5\}$.
\end{solution}
\end{question}

\begin{question}[points=4]
已知$f(x)=\begin{cases} \sqrt{x}, & x>0 \\ 1, & x\le 0 \end{cases}$，则$f(3)=$.
\fillin[{$3$}]
\begin{solution}
\textbf{答案：}$3$

\textbf{分析：}利用分段函数的形式可求$f(3)$.

\textbf{详解：}因为$f(x)=\begin{cases} \sqrt{x}, & x>0 \\ 1, & x\le 0 \end{cases}$，故$f(3)=\sqrt{3}=3$，故答案为：$3$.
\end{solution}
\end{question}

\begin{question}[points=4]
已知$x\in\mathbb{R}$，则不等式$x^2-2x-3<0$的解集为.
\fillin[{$\{x\mid -1<x<3\}$}]
\begin{solution}
\textbf{答案：}$\{x\mid -1<x<3\}$

\textbf{分析：}求出方程$x^2-2x-3=0$的解后可求不等式的解集.

\textbf{详解：}方程$x^2-2x-3=0$的解为$x=-1$或$x=3$，故不等式$x^2-2x-3<0$的解集为$\{x\mid -1<x<3\}$，故答案为：$\{x\mid -1<x<3\}$.
\end{solution}
\end{question}

\begin{question}[points=4]
已知$f(x)=x^3+a$，$x\in\mathbb{R}$，且$f(x)$是奇函数，则$a=$.
\fillin[{$0$}]
\begin{solution}
\textbf{答案：}$0$

\textbf{分析：}根据奇函数的性质可求参数$a$.

\textbf{详解：}因为$f(x)$是奇函数，故$f(-x)+f(x)=0$即$(-x)^3+a+x^3+a=0$，故$a=0$，故答案为：$0$.
\end{solution}
\end{question}

\begin{question}[points=4]
已知$k\in\mathbb{R}$，$\vec{a}=(2,5)$，$\vec{b}=(6,k)$，且$\vec{a}\parallel\vec{b}$，则$k$的值为.
\fillin[{$15$}]
\begin{solution}
\textbf{答案：}$15$

\textbf{分析：}根据向量平行的坐标表示得到方程，解出即可.

\textbf{详解：}$\because\vec{a}\parallel\vec{b}$，$\therefore 2k=5\times6$，解得$k=15$．故答案为：15.
\end{solution}
\end{question}

\begin{question}[points=4]
在$(x+1)^n$的二项展开式中，若各项系数和为32，则$x^2$项的系数为.
\fillin[{$10$}]
\begin{solution}
\textbf{答案：}$10$

\textbf{分析：}令$x=1$，解出$n=5$，再利用二项式的展开式的通项合理赋值即可.

\textbf{详解：}令$x=1$，$\therefore (1+1)^n=32$，即$2^n=32$，解得$n=5$，所以$(x+1)^5$的展开式通项公式为$T_{r+1}=C_5^r x^{5-r}$，令$5-r=2$，则$r=3$，$\therefore T_4=C_5^3 x^2=10x^2$．故答案为：10.
\end{solution}
\end{question}

\begin{question}[points=5]
已知抛物线$y^2=4x$上有一点$P$到准线的距离为9，那么点$P$到$x$轴的距离为.
\fillin[{$4\sqrt{2}$}]
\begin{solution}
\textbf{答案：}$4\sqrt{2}$

\textbf{分析：}根据抛物线的定义知$x_P=8$，将其再代入抛物线方程即可.

\textbf{详解：}由$y^2=4x$知抛物线的准线方程为$x=-1$，设点$P(x_0,y_0)$，由题意得$x_0+1=9$，解得$x_0=8$，代入抛物线方程$y^2=4x$，得$y_0^2=32$，解得$y_0=\pm4\sqrt{2}$，则点$P$到$x$轴的距离为$4\sqrt{2}$．故答案为：$4\sqrt{2}$.
\end{solution}
\end{question}

\begin{question}[points=5]
某校举办科学竞技比赛，有A、B、C三种题库，A题库有5000道题，B题库有4000道题，C题库有3000道题．小申已完成所有题，他A题库的正确率是0.92，B题库的正确率是0.86，C题库的正确率是0.72．现他从所有的题中随机选一题，正确率是.
\fillin[{$0.85$}]
\begin{solution}
\textbf{答案：}$0.85$

\textbf{分析：}求出各题库所占比，根据全概率公式即可得到答案.

\textbf{详解：}由题意知，A,B,C题库的比例为：$5:4:3$，各占比分别为$\frac{5}{12},\frac{4}{12},\frac{3}{12}$，则根据全概率公式知所求正确率$p=\frac{5}{12}\times0.92+\frac{4}{12}\times0.86+\frac{3}{12}\times0.72=0.85$．故答案为：0.85.
\end{solution}
\end{question}

\begin{question}[points=5]
已知虚数$z$，其实部为1，且$z+\frac{2}{z}=m(m\in\mathbb{R})$，则实数$m$为.
\fillin[{$2$}]
\begin{solution}
\textbf{答案：}$2$

\textbf{分析：}设$z=1+bi$，直接根据复数的除法运算，再根据复数分类即可得到答案.

\textbf{详解：}设$z=1+bi$，$b\in\mathbb{R}$且$b\neq0$．则$z+\frac{2}{z}=1+bi+\frac{2}{1+bi}=\left(\frac{b^2+3}{1+b^2}\right)+\left(\frac{b^3-b}{1+b^2}\right)i=m$，$\because m\in\mathbb{R}$，$\therefore \begin{cases} \frac{b^2+3}{1+b^2}=m \\ \frac{b^3-b}{1+b^2}=0 \end{cases}$，解得$m=2$，故答案为：2.
\end{solution}
\end{question}

\begin{question}[points=5]
设集合$A$中的元素皆为无重复数字的三位正整数，且元素中任意两者之积皆为偶数，求集合中元素个数的最大值.
\fillin[{$329$}]
\begin{solution}
\textbf{答案：}$329$

\textbf{分析：}三位数中的偶数分个位是0和个位不是0讨论即可.

\textbf{详解：}由题意知集合中且至多只有一个奇数，其余均是偶数．首先讨论三位数中的偶数，①当个位为0时，则百位和十位在剩余的9个数字中选择两个进行排列，则这样的偶数有$P_9^2=72$个；②当个位不为0时，则个位有$C_4^1$个数字可选，百位有$C_8^1$个数字可选，十位有$C_8^1$个数字可选，根据分步乘法这样的偶数共有$C_4^1C_8^1C_8^1=256$，最后再加上单独的奇数，所以集合中元素个数的最大值为$72+256+1=329$个．故答案为：329.
\end{solution}
\end{question}

\begin{question}[points=5]
已知点$B$在点$C$正北方向，点$D$在点$C$的正东方向，$BC=CD$，存在点$A$满足$\angle BAC=16.5^\circ$，$\angle DAC=37^\circ$，则$\angle BCA=$(精确到0.1度).
\fillin[{$7.8^\circ$}]
\begin{solution}
\textbf{答案：}$7.8^\circ$

\textbf{分析：}设$\angle BCA=\theta$，在$\triangle DCA$和$\triangle BCA$中分别利用正弦定理得到$\frac{CA}{\sin D}=\frac{CD}{\sin\angle CAD}$，$\frac{CA}{\sin(\theta+16.5^\circ)}=\frac{CB}{\sin16.5^\circ}$，两式相除即可得到答案.

\textbf{详解：}设$\angle BCA=\theta$，$\angle ACD=90^\circ-\theta$，在$\triangle DCA$中，由正弦定理得$\frac{CA}{\sin D}=\frac{CD}{\sin\angle CAD}$，即$\frac{CA}{\sin[180^\circ-(90^\circ-\theta+37.0^\circ)]}=\frac{CD}{\sin37.0^\circ}$，即$\frac{CA}{\sin(90^\circ-\theta+37.0^\circ)}=\frac{CD}{\sin37.0^\circ}$①；在$\triangle BCA$中，由正弦定理得$\frac{CA}{\sin B}=\frac{CB}{\sin\angle CAB}$，即$\frac{CA}{\sin[180^\circ-(\theta+16.5^\circ)]}=\frac{CB}{\sin16.5^\circ}$，即$\frac{CA}{\sin(\theta+16.5^\circ)}=\frac{CB}{\sin16.5^\circ}$②，因为$CD=CB$，$\frac{②}{①}$得$\frac{\sin(90^\circ-\theta+37.0^\circ)}{\sin(\theta+16.5^\circ)}=\frac{\sin37.0^\circ}{\sin16.5^\circ}$，利用计算器即可得$\theta\approx7.8^\circ$，故答案为：$7.8^\circ$.
\end{solution}
\end{question}

\begin{question}[points=5]
无穷等比数列$\{a_n\}$满足首项$a_1>0$，$q>1$，记$I_n=\{x-y\mid x,y\in[a_1,a_2]\cup[a_n,a_{n+1}]\}$，若对任意正整数$n$集合$I_n$是闭区间，则$q$的取值范围是.
\fillin[{$q\ge2$}]
\begin{solution}
\textbf{答案：}$q\ge2$

\textbf{分析：}当$n\ge2$时，不妨设$x\ge y$，则$x-y\in[0,a_2-a_1]\cup[a_n-a_2,a_{n+1}-a_1]\cup[0,a_{n+1}-a_n]$，结合$I_n$为闭区间可得$q-2\ge-\frac{1}{q^{n-2}}$对任意的$n\ge2$恒成立，故可求$q$的取值范围.

\textbf{详解：}由题设有$a_n=a_1q^{n-1}$，因为$a_1>0,q>1$，故$a_{n+1}>a_n$，故$[a_n,a_{n+1}]=[a_1q^{n-1},a_1q^n]$，当$n=1$时，$x,y\in[a_1,a_2]$，故$x-y\in[a_1-a_2,a_2-a_1]$，此时$I_1$为闭区间，当$n\ge2$时，不妨设$x\ge y$，若$x,y\in[a_1,a_2]$，则$x-y\in[0,a_2-a_1]$，若$y\in[a_1,a_2],x\in[a_n,a_{n+1}]$，则$x-y\in[a_n-a_2,a_{n+1}-a_1]$，若$x,y\in[a_n,a_{n+1}]$，则$x-y\in[0,a_{n+1}-a_n]$，综上，$x-y\in[0,a_2-a_1]\cup[a_n-a_2,a_{n+1}-a_1]\cup[0,a_{n+1}-a_n]$，又$I_n$为闭区间等价于$[0,a_2-a_1]\cup[a_n-a_2,a_{n+1}-a_1]\cup[0,a_{n+1}-a_n]$为闭区间，而$a_{n+1}-a_1>a_{n+1}-a_n>a_2-a_1$，故$a_{n+1}-a_n\ge a_n-a_2$对任意$n\ge2$恒成立，故$a_{n+1}-2a_n+a_2\ge0$即$a_1q^{n-1}(q-2)+a_2\ge0$，故$q^{n-2}(q-2)+1\ge0$，故$q-2\ge-\frac{1}{q^{n-2}}$对任意的$n\ge2$恒成立，因$q>1$，故当$n\to+\infty$时，$-\frac{1}{q^{n-2}}\to0$，故$q-2\ge0$即$q\ge2$．故答案为：$q\ge2$.
\end{solution}
\end{question}

\section*{二、选择题}
本大题共有4题，满分18分，其中第13-14题每题4分，第15-16题每题5分。每题有且只有一个正确答案，考生应在答题纸的相应编号上，将代表答案的小方格涂黑，选对得满分，否则一律得零分.

\begin{question}[points=4]
已知气候温度和海水表层温度相关，且相关系数为正数，对此描述正确的是（\quad\quad）
\begin{choices}
  \item $A$气候温度高，海水表层温度就高
  \item $B$气候温度高，海水表层温度就低
  \item $C$随着气候温度由低到高，海水表层温度呈上升趋势
  \item $D$随着气候温度由低到高，海水表层温度呈下降趋势
\end{choices}
\begin{solution}
\textbf{答案：}$C$

\textbf{分析：}根据相关系数的性质可得正确的选项.

\textbf{详解：}对于AB，当气候温度高，海水表层温度变高变低不确定，故AB错误.对于CD，因为相关系数为正，故随着气候温度由低到高时，海水表层温度呈上升趋势，故C正确，D错误.故选：C.
\end{solution}
\end{question}

\begin{question}[points=4]
下列函数$f(x)$的最小正周期是$2\pi$的是（\quad\quad）
\begin{choices}
  \item $A\ \sin x+\cos x$
  \item $B\ \sin x\cos x$
  \item $C\ \sin^2x+\cos^2x$
  \item $D\ \sin^2x-\cos^2x$
\end{choices}
\begin{solution}
\textbf{答案：}$A$

\textbf{分析：}根据辅助角公式、二倍角公式以及同角三角函数关系并结合三角函数的性质一一判断即可.

\textbf{详解：}对A，$\sin x+\cos x=\sqrt{2}\sin\left(x+\frac{\pi}{4}\right)$，周期$T=2\pi$，故A正确；对B，$\sin x\cos x=\frac{1}{2}\sin2x$，周期$T=\frac{2\pi}{2}=\pi$，故B错误；对于选项C，$\sin^2x+\cos^2x=1$，是常值函数，不存在最小正周期，故C错误；对于选项D，$\sin^2x-\cos^2x=-\cos2x$，周期$T=\frac{2\pi}{2}=\pi$，故D错误，故选：A.
\end{solution}
\end{question}

\begin{question}[points=5]
定义一个集合$\Omega$，集合中的元素是空间内的点集，任取$P_1,P_2,P_3\in\Omega$，存在不全为0的实数$\lambda_1,\lambda_2,\lambda_3$，使得$\lambda_1\overrightarrow{OP_1}+\lambda_2\overrightarrow{OP_2}+\lambda_3\overrightarrow{OP_3}=\vec{0}$．已知$(1,0,0)\in\Omega$，则$(0,0,1)\notin\Omega$的充分条件是（\quad\quad）
\begin{choices}
  \item $A\ (0,0,0)\in\Omega$
  \item $B\ (-1,0,0)\in\Omega$
  \item $C\ (0,1,0)\in\Omega$
  \item $D\ (0,0,-1)\in\Omega$
\end{choices}
\begin{solution}
\textbf{答案：}$C$

\textbf{分析：}首先分析出三个向量共面，显然当$(1,0,0),(0,0,1),(0,1,0)\in\Omega$时，三个向量构成空间的一个基底，则即可分析出正确答案.

\textbf{详解：}由题意知这三个向量$\overrightarrow{OP_1},\overrightarrow{OP_2},\overrightarrow{OP_3}$共面，即这三个向量不能构成空间的一个基底，对A，由空间直角坐标系易知$(0,0,0),(1,0,0),(0,0,1)$三个向量共面，则当$(-1,0,0),(1,0,0)\in\Omega$无法推出$(0,0,1)\notin\Omega$，故A错误；对B，由空间直角坐标系易知$(-1,0,0),(1,0,0),(0,0,1)$三个向量共面，则当$(0,0,0),(1,0,0)\in\Omega$无法推出$(0,0,1)\notin\Omega$，故A错误；对C，由空间直角坐标系易知$(1,0,0),(0,0,1),(0,1,0)$三个向量不共面，可构成空间的一个基底，则由$(1,0,0),(0,1,0)\in\Omega$能推出$(0,0,1)\notin\Omega$，对D，由空间直角坐标系易知$(1,0,0),(0,0,1),(0,0,-1)$三个向量共面，则当$(0,0,-1)(1,0,0)\in\Omega$无法推出$(0,0,1)\notin\Omega$，故D错误.故选：C.
\end{solution}
\end{question}

\begin{question}[points=5]
已知函数$f(x)$的定义域为$\mathbb{R}$，定义集合$M=\{x_0\mid x_0\in\mathbb{R},\forall x\in(-\infty,x_0),f(x)<f(x_0)\}$，在使得$M=[-1,1]$的所有$f(x)$中，下列成立的是（\quad\quad）
\begin{choices}
  \item $A$存在$f(x)$是偶函数
  \item $B$存在$f(x)$在$x=2$处取最大值
  \item $C$存在$f(x)$是严格增函数
  \item $D$存在$f(x)$在$x=-1$处取到极小值
\end{choices}
\begin{solution}
\textbf{答案：}$B$

\textbf{分析：}对于ACD利用反证法并结合函数奇偶性、单调性以及极小值的概念即可判断，对于B，构造函数$f(x)=\begin{cases} -2, & x<-1 \\ x, & -1\le x\le1 \\ 1, & x>1 \end{cases}$即可判断.

\textbf{详解：}对于A，若存在$y=f(x)$是偶函数，取$x_0=1\in[-1,1]$，则对于任意$x\in(-\infty,1),f(x)<f(1)$，而$f(-1)=f(1)$，矛盾，故A错误；对于B，可构造函数$f(x)=\begin{cases} -2, & x<-1 \\ x, & -1\le x\le1 \\ 1, & x>1 \end{cases}$，满足集合$M=[-1,1]$，当$x<-1$时，则$f(x)=-2$，当$-1\le x\le1$时，$f(x)\in[-1,1]$，当$x>1$时，$f(x)=1$，则该函数$f(x)$的最大值是$f(2)$，则B正确；对C，假设存在$f(x)$，使得$f(x)$严格递增，则$M=\mathbb{R}$，与已知$M=[-1,1]$矛盾，则C错误；对D，假设存在$f(x)$，使得$f(x)$在$x=-1$处取极小值，则在-1的左侧附近存在$n$，使得$f(n)>f(-1)$，这与已知集合$M$的定义矛盾，故D错误；故选：B.
\end{solution}
\end{question}

\section*{三、解答题}
本大题共有5题，满分78分。解下列各题必须在答题纸相应编号的规定区域内写出必要的步骤.

\begin{question}[points=14]
如图为正四棱锥$P-ABCD$，$O$为底面$ABCD$的中心．

\begin{center}
  \includegraphics[width=.38\linewidth]{question-17-1.jpg}
\end{center}
\begin{solution}
\textbf{答案：}（1）$12\pi$；（2）$\frac{\pi}{4}$

\textbf{分析：}（1）根据正四棱锥的数据，先算出直角三角形$\triangle POA$的边长，然后求圆锥的体积；（2）连接$EA,EO,EC$，可先证$BE\perp$平面$ACE$，根据线面角的定义得出所求角为$\angle BOE$，然后结合题目数量关系求解.

\textbf{详解：}（1）正四棱锥满足且$PO\perp$平面$ABCD$，由$AO\subset$平面$ABCD$，则$PO\perp AO$，又正四棱锥底面$ABCD$是正方形，由$AD=3\sqrt2$可得，$AO=3$，故$PO=\sqrt{PA^2-AO^2}=4$，根据圆锥的定义，$\triangle POA$绕$PO$旋转一周形成的几何体是以$PO$为轴，$AO$为底面半径的圆锥，即圆锥的高为$PO=4$，底面半径为$AO=3$，根据圆锥的体积公式，所得圆锥的体积是$\frac13\times\pi\times3^2\times4=12\pi$.（2）连接$EA,EO,EC$，由题意结合正四棱锥的性质可知，每个侧面都是等边三角形，由$E$是$PB$中点，则$AE\perp PB$，$CE\perp PB$，又$AE\cap CE=E$，$AE,CE\subset$平面$ACE$，故$PB\perp$平面$ACE$，即$BE\perp$平面$ACE$，又$BD\cap$平面$ACE=O$，于是直线$BD$与平面$AEC$所成角的大小即为$\angle BOE$，不妨设$AP=AD=6$，则$BO=3\sqrt2$，$BE=3$，$\sin\angle BOE=\frac{3}{3\sqrt2}=\frac{\sqrt2}{2}$，又线面角的范围是$[0,\frac{\pi}{2}]$，故$\angle BOE=\frac{\pi}{4}$，即为所求.
\end{solution}
\end{question}

\begin{question}[points=14]
若$f(x)=\log_a x$（$a>0$a≠1）．
\begin{solution}
\textbf{答案：}（1）$\{x\mid1<x<2\}$；（2）$a>1$

\textbf{分析：}（1）求出底数$a$，再根据对数函数的单调性可求不等式的解；（2）存在$x$使得$f(x+1)$，$f(ax)$，$f(x+2)$成等差数列等价于$a^2=2\left(\frac1x+\frac34\right)^2-\frac18$在$(0,+\infty)$上有解，利用换元法结合二次函数的性质可求$a$的取值范围.

\textbf{详解：}（1）因为$y=f(x)$的图象过$(4,2)$，故$\log_a4=2$，故$a^2=4$即$a=2$（负的舍去），而$f(x)=\log_2x$在$(0,+\infty)$上为增函数，故$f(2x-2)<f(x)$，故$0<2x-2<x$即$1<x<2$，故$f(2x-2)<f(x)$的解集为$\{x\mid1<x<2\}$.（2）因为存在$x$使得$f(x+1)$，$f(ax)$，$f(x+2)$成等差数列，故$2f(ax)=f(x+1)+f(x+2)$有解，故$2\log_a(ax)=\log_a(x+1)+\log_a(x+2)$，因为$a>0$，$a\neq1$，故$x>0$，故$a^2x^2=(x+1)(x+2)$在$(0,+\infty)$上有解，由$a^2=\frac{x^2+3x+2}{x^2}=1+\frac3x+\frac2{x^2}=2\left(\frac1x+\frac34\right)^2-\frac18$在$(0,+\infty)$上有解，令$t=\frac1x\in(0,+\infty)$，而$y=2\left(t+\frac34\right)^2-\frac18$在$(0,+\infty)$上的值域为$(1,+\infty)$，故$a^2>1$即$a>1$.
\end{solution}
\end{question}

\begin{question}[points=14]
为了解某地初中学生体育锻炼时长与学业成绩的关系，从该地区29000名学生中抽取580人，得到日均体育锻炼时长与学业成绩的数据如下表所示：\begin{center}\begin{tabular}{c|ccccc} \text{时间范围} & [0,0.5) & [0.5,1) & [1,1.5) & [1.5,2) & [2,2.5] \\ \hline \text{优秀} & 5 & 44 & 42 & 3 & 1 \\ \text{不优秀} & 134 & 147 & 137 & 40 & 27 \end{tabular}\end{center}（1）该地区29000名学生中体育锻炼时长不少于1小时人数约为多少？（2）估计该地区初中学生日均体育锻炼的时长（精确到0.1）；（3）是否有95%的把握认为学业成绩优秀与日均体育锻炼时长不小于1小时且小于2小时有关？附：$\chi^2=\frac{n(ad-bc)^2}{(a+b)(c+d)(a+c)(b+d)}$，其中$n=a+b+c+d$，$P(\chi^2\ge3.841)\approx0.05$.
\begin{solution}
\textbf{答案：}（1）12500；（2）0.9h；（3）有

\textbf{分析：}（1）求出相关占比，乘以总人数即可；（2）根据平均数的计算公式即可得到答案；（3）作出列联表，再提出零假设，计算卡方值和临界值比较大小即可得到结论.

\textbf{详解：}（1）由表可知锻炼时长不少于1小时的人数为占比$\frac{179+43+28}{580}=\frac{25}{58}$，则估计该地区29000名学生中体育锻炼时长不少于1小时的人数为$29000\times\frac{25}{58}=12500$．（2）估计该地区初中生的日均体育锻炼时长约为$\frac1{580}[\frac{0.5}{2}\times139+\frac{0.5+1}{2}\times191+\frac{1+1.5}{2}\times179+\frac{1.5+2}{2}\times43+\frac{2+2.5}{2}\times28]\approx0.9$．则估计该地区初中学生日均体育锻炼的时长为0.9小时.（3）由题列联表如下：\begin{center}\begin{tabular}{c|cc|c} & [1,2) & 其他 & 合计 \\ \hline 优秀 & 45 & 50 & 95 \\ 不优秀 & 177 & 308 & 485 \\ \hline 合计 & 222 & 358 & 580 \end{tabular}\end{center}提出零假设$H_0$：该地区成绩优秀与日均锻炼时长不少于1小时但少于2小时无关．其中$\alpha=0.05$．$\chi^2=\frac{580\times(45\times308-177\times50)^2}{95\times485\times222\times358}\approx3.976>3.841$．则零假设不成立，即有95%的把握认为学业成绩优秀与日均锻炼时长不小于1小时且小于2小时有关．
\end{solution}
\end{question}

\begin{question}[points=18]
已知双曲线$\Gamma:x^2-\frac{y^2}{b^2}=1$，$(b>0)$，左右顶点分别为$A_1$，$A_2$，过点$M(-2,0)$的直线$l$交双曲线$\Gamma$于$P$，$Q$两点.（1）若离心率$e=2$时，求$b$的值．（2）若$b=\frac{2\sqrt6}{3}$，$\triangle MA_2P$为等腰三角形时，且点$P$在第一象限，求点$P$的坐标．（3）连接$OQ$并延长，交双曲线$\Gamma$于点$R$，若$\overrightarrow{A_1R}\cdot\overrightarrow{A_2P}=1$，求$b$的取值范围．
\begin{solution}
\textbf{答案：}（1）$b=\sqrt3$；（2）$P(2,2\sqrt2)$；（3）$(0,\sqrt3)\cup(\sqrt3,\frac{\sqrt{30}}{3}]$

\textbf{分析：}（1）根据离心率公式计算即可；（2）分三角形三边分别为底讨论即可；（3）设直线$l:x=my-2$，联立双曲线方程得到韦达定理式，再代入计算向量数量积的等式计算即可.

\textbf{详解：}（1）由题意得$e=\frac{c}{a}=c=2$，则$c=2$，$b=\sqrt{2^2-1}=\sqrt3$．（2）当$b=\frac{2\sqrt6}{3}$时，双曲线$\Gamma:x^2-\frac{y^2}{\frac{8}{3}}=1$，其中$M(-2,0)$，$A_2(1,0)$，因为$\triangle MA_2P$为等腰三角形，则①当以$MA_2$为底时，显然点$P$在直线$x=-\frac12$上，这与点$P$在第一象限矛盾，故舍去；②当以$A_2P$为底时，$MP=MA_2=3$，设$P(x,y)$，则$\begin{cases} x^2-\frac{3y^2}{8}=1 \\ (x+2)^2+y^2=9 \end{cases}$，联立解得$\begin{cases} x=-\frac{23}{11} \\ y=-\frac{8\sqrt{17}}{11} \end{cases}$或$\begin{cases} x=-\frac{23}{11} \\ y=\frac{8\sqrt{17}}{11} \end{cases}$或$\begin{cases} x=1 \\ y=0 \end{cases}$，因为点$P$在第一象限，显然以上均不合题意，舍去；（或者由双曲线性质知$MP>MA_2$，矛盾，舍去）；③当以$MP$为底时，$A_2P=MA_2=3$，设$P(x_0,y_0)$，其中$x_0>0,y_0>0$，则有$\begin{cases} (x_0-1)^2+y_0^2=9 \\ x_0^2-\frac{3y_0^2}{8}=1 \end{cases}$，解得$\begin{cases} x_0=2 \\ y_0=2\sqrt2 \end{cases}$，即$P(2,2\sqrt2)$．综上所述：$P(2,2\sqrt2)$.（3）由题知$A_1(-1,0)$，$A_2(1,0)$，当直线$l$的斜率为0时，此时$\overrightarrow{A_1R}\cdot\overrightarrow{A_2P}=0$，不合题意，则$k_l\neq0$，则设直线$l:x=my-2$，设点$P(x_1,y_1)$，$Q(x_2,y_2)$，根据$OQ$延长线交双曲线$\Gamma$于点$R$，根据双曲线对称性知$R(-x_2,-y_2)$，联立有$\begin{cases} x=my-2 \\ x^2-\frac{y^2}{b^2}=1 \end{cases}\Rightarrow (b^2m^2-1)y^2-4b^2my+3b^2=0$，显然二次项系数$b^2m^2-1\neq0$，其中$\Delta=(-4mb^2)^2-4(b^2m^2-1)3b^2=4b^4m^2+12b^2>0$，$y_1+y_2=\frac{4b^2m}{b^2m^2-1}①$，$y_1y_2=\frac{3b^2}{b^2m^2-1}②$，$\overrightarrow{A_1R}=(-x_2+1,-y_2)$，$\overrightarrow{A_2P}=(x_1-1,y_1)$，则$\overrightarrow{A_1R}\cdot\overrightarrow{A_2P}=(-x_2+1)(x_1-1)-y_1y_2=1$，因为$P(x_1,y_1)$，$Q(x_2,y_2)$在直线$l$上，则$x_1=my_1-2$，$x_2=my_2-2$，即$-(my_2-3)(my_1-3)-y_1y_2=1$，即$y_1y_2(m^2+1)-3m(y_1+y_2)+10=0$，将①②代入有$(m^2+1)\frac{3b^2}{b^2m^2-1}-3m\frac{4b^2m}{b^2m^2-1}+10=0$，即$3b^2(m^2+1)-12b^2m^2+10(b^2m^2-1)=0$，化简得$b^2m^2+3b^2-10=0$，所以$m^2=\frac{10}{b^2}-3$，代入到$b^2m^2-1\neq0$，得$b^2(\frac{10}{b^2}-3)-1=10-3b^2-1=9-3b^2\neq0$，所以$b^2\neq3$，且$m^2=\frac{10}{b^2}-3\ge0$，解得$b^2\le\frac{10}{3}$，又因为$b>0$，则$0<b\le\frac{\sqrt{30}}{3}$，综上知，$b\in(0,\sqrt3)\cup(\sqrt3,\frac{\sqrt{30}}{3}]$.
\end{solution}
\end{question}

\begin{question}[points=18]
对于一个函数$f(x)$和一个点$M(a,b)$，令$s(x)=(x-a)^2+(f(x)-b)^2$，若$P(x_0,f(x_0))$是$s(x)$取到最小值的点，则称$P$是$M$在$f(x)$的“最近点”．（1）对于$f(x)=\frac1x$（$x>0$），求证：对于点$M(0,0)$，存在点$P$，使得点$P$是$M$在$f(x)$的“最近点”；（2）对于$f(x)=e^x$，$M(1,0)$，请判断是否存在一个点$P$，它是$M$在$f(x)$的“最近点”，且直线$MP$与$y=f(x)$在点$P$处的切线垂直；（3）已知$y=f(x)$在定义域$\mathbb{R}$上存在导函数$f'(x)$，且函数$g(x)$在定义域$\mathbb{R}$上恒正，设点$M_1(t-1,f(t)-g(t))$，$M_2(t+1,f(t)+g(t))$．若对任意的$t\in\mathbb{R}$，存在点$P$同时是$M_1$，$M_2$在$f(x)$的“最近点”，试判断$f(x)$的单调性．
\begin{solution}
\textbf{答案：}（1）证明见解析；（2）存在，$P(0,1)$；（3）严格单调递减

\textbf{分析：}（1）代入M(0,0)，利用基本不等式即可；（2）由题得$s(x)=(x-1)^2+e^{2x}$，利用导函数得到其最小值，则得到P，再证明直线MP与切线垂直即可；（3）根据题意得到$s_1'(x_0)=s_2'(x_0)=0$，对两等式化简得$f'(x_0)=-\frac1{g(t)}$，再利用“最近点”的定义得到不等式组，即可证明$x_0=t$，最后得到函数单调性.

\textbf{详解：}（1）当M(0,0)时，$s(x)=(x-0)^2+(\frac1x-0)^2=x^2+\frac1{x^2}\ge2\sqrt{x^2\cdot\frac1{x^2}}=2$，当且仅当$x^2=\frac1{x^2}$即$x=1$时取等号，故对于点M(0,0)，存在点P(1,1)，使得该点是M(0,0)在$f(x)$的“最近点”．（2）由题设可得$s(x)=(x-1)^2+(e^x-0)^2=(x-1)^2+e^{2x}$，则$s'(x)=2(x-1)+2e^{2x}$，因为$y=2(x-1)$，$y=2e^{2x}$均为R上单调递增函数，则$s'(x)=2(x-1)+2e^{2x}$在R上为严格增函数，而$s'(0)=0$，故当$x<0$时，$s'(x)<0$，当$x>0$时，$s'(x)>0$，故$s(x)_{\min}=s(0)=2$，此时P(0,1)，而$f'(x)=e^x$，$k=f'(0)=1$，故$f(x)$在点P处的切线方程为$y=x+1$．而$k_{MP}=\frac{0-1}{1-0}=-1$，故$k_{MP}\cdot k=-1$，故直线MP与$y=f(x)$在点P处的切线垂直．（3）设$s_1(x)=(x-t+1)^2+(f(x)-f(t)+g(t))^2$，$s_2(x)=(x-t-1)^2+(f(x)-f(t)-g(t))^2$，而$s_1'(x)=2(x-t+1)+2(f(x)-f(t)+g(t))f'(x)$，$s_2'(x)=2(x-t-1)+2(f(x)-f(t)-g(t))f'(x)$，若对任意的$t\in\mathbb{R}$，存在点P同时是$M_1$，$M_2$在$f(x)$的“最近点”，设$P(x_0,y_0)$，则$x_0$既是$s_1(x)$的最小值点，也是$s_2(x)$的最小值点，因为两函数的定义域均为R，则$x_0$也是两函数的极小值点，则存在$x_0$，使得$s_1'(x_0)=s_2'(x_0)=0$，即$s_1'(x_0)=2(x_0-t+1)+2f'(x_0)[f(x_0)-f(t)+g(t)]=0$①，$s_2'(x_0)=2(x_0-t-1)+2f'(x_0)[f(x_0)-f(t)-g(t)]=0$②，由①②相等得$4+4g(t)f'(x_0)=0$，即$1+f'(x_0)g(t)=0$，即$f'(x_0)=-\frac1{g(t)}$，又因为函数$g(x)$在定义域R上恒正，则$f'(x_0)=-\frac1{g(t)}<0$恒成立，接下来证明$x_0=t$，因为$x_0$既是$s_1(x)$的最小值点，也是$s_2(x)$的最小值点，则$s_1(x_0)\le s_1(t)$，$s_2(x_0)\le s_2(t)$，即$(x_0-t+1)^2+(f(x_0)-f(t)+g(t))^2\le1+(g(t))^2$③，$(x_0-t-1)^2+(f(x_0)-f(t)-g(t))^2\le1+(g(t))^2$④，③+④得$2(x_0-t)^2+2+2[f(x_0)-f(t)]^2+2g^2(t)\le2+2g^2(t)$，即$(x_0-t)^2+[f(x_0)-f(t)]^2\le0$，因为$(x_0-t)^2\ge0$，$[f(x_0)-f(t)]^2\ge0$，则$\begin{cases} x_0-t=0 \\ f(x_0)-f(t)=0 \end{cases}$，解得$x_0=t$，则$f'(t)=-\frac1{g(t)}<0$恒成立，因为t的任意性，则$f(x)$严格单调递减.
\end{solution}
\end{question}

